\documentclass{article}

\title{Java III Test II Instructions}
\author{Scott R. Young}
\date{\today}

\begin{document}

\maketitle

You are limited to using techniques and tools taught in this course (and its
prerequisites) through lectures, examples, and evaluations.

Prompt driven Generative AIs, such as Co-pilot and Claude, are
strictly prohibited for this test.

You are only to modify code that is marked with TODO comments.

With a built .war file deployed on a TomEE 10 server, the test should print out
what steps it can complete and any exceptions it encounters when you generate
user events on the UI.

\section{Instructions}
\begin{enumerate}
    \item Download and unzip the test2.zip file for this test.
    \item Open the folder containing db.sql in VS Code.
    \item Open HeidiSQL and open your MariaDB connection.
    \item Open a query in HeidiSQL.
    \item Copy the included db.sql file's contents into the query.
    \item Execute the query.
    \item Open the following file in VS Code: \\
            \textbf{test2/src/main/webapp/WEB-INF/web.xml}.
    \item Follow the directions there to allow your application and service to connect to your DB.
    \item Implement the TODO comments in:\\
            \textbf{test2/src/main/webapp/index.xhtml}\\
            \textbf{test2/src/main/java/test2/app/ShoppingCart.java}\\
            \textbf{test2/src/main/java/test2/app/User.java}\\
            \textbf{test2/src/main/java/test2/service/CartRootResource.java}\\
            \textbf{test2/src/main/java/test2/service/ItemRootResource.java}
\end{enumerate}

\section{Additional Information}
This test is open book; open notes.

You may use anything on the D2L shell, the Java SE docs, the Jakarta EE docs, the official Jakarta EE tutorials (https:
//jakarta.ee/learn/docs/jakartaee-tutorial/current/index.html), or any notes you've written yourself ahead of time.

\end{document}